\chapter{System Design}
\label{ch:design}

This chapter describes A.L.F.R.E.D.\ as a system: the two execution tiers, the
artifacts that connect them, and the design principles that hold the whole thing
together. The goal is to fix the vocabulary that the methodology and evaluation
chapters reuse, and to make one architectural claim concrete before any number
is reported: \emph{most on-device requests do not need a model to reason; they
need a model to fill in a blank}.

\section{Overview: two tiers, one router}
\label{sec:overview}

A.L.F.R.E.D.\ sits between a user's natural-language request and the device's
capabilities. Every request enters a \textbf{router}, which makes a single
decision: can this request be served by a known, repeatable \emph{pattern}, or
does it require open-ended problem solving? The first case is dispatched to the
\textbf{reflexer}; the second to the \textbf{thinker}.

\begin{figure}[t]
  \centering
  \includegraphics[width=0.95\textwidth]{fig_architecture.pdf}
  \caption{The two-tier architecture. The router classifies each request; the
  reflexer handles the deterministic head with a distilled pattern, the thinker
  handles the long tail with a full agent loop. Both emit commands to the same
  skill CLIs.}
  \label{fig:architecture}
\end{figure}

The split is an economic bet, not a capability ranking. The reflexer is cheap,
fast, and runs comfortably on the device; it is also narrow, because it can only
serve intents for which a pattern exists. The thinker is general but expensive:
it runs an agent loop, costs an order of magnitude more tokens and latency per
request, and---at the sizes that reason well---does not fit on the device at
all. Neither tier is sufficient alone. The contribution of the architecture is
the \emph{division of labour}: route the high-frequency, repeatable head to the
reflexer and reserve the thinker for the tail. Chapter~\ref{ch:eval} quantifies
exactly where that bet pays.

\section{The reflexer: an executioner, not a generator}
\label{sec:reflexer}

The reflexer is a small language model (in this work, \texttt{qwen-3.5-2b})
whose job is deliberately impoverished. It does not decide \emph{what} to do or
\emph{how} to do it. Given a request that the router has already matched to a
pattern, the reflexer extracts the variable parts of the request---the
\emph{slots}---and substitutes them into a fixed command template supplied by
that pattern. It then emits exactly one shell command. This is why we call it an
\emph{executioner} rather than a generator: it carries out a decision that was
made once, offline, by something else.

This framing is the central design choice of the thesis, and it separates
A.L.F.R.E.D.\ from the on-device assistant models shipped by Google and Apple.
Those models are asked to be creative within a constrained surface; they
generate. The reflexer is asked for none of that. It performs slot extraction
and template binding, and nothing more. The benefit is reliability: the space of
things the reflexer can get wrong shrinks to the space of slot values, because
the command shape itself is not the model's to invent. The cost is generality,
paid up front as the requirement that a pattern must exist
(Section~\ref{sec:distillation}).

Concretely, the runtime is the short pipeline in
Figure~\ref{fig:reflexer}. The matched pattern supplies a \emph{command shape}
(the \texttt{args\_template} with its slots rendered as typed placeholders) and a
list of slot descriptions; these, with the user query, form the entire prompt to
the 2B model. The model emits one shell command directly---there is no JSON step
and no separate template-rendering pass. The command is executed once through a
shell (with a per-step timeout and an optional retry), and the pattern's
\emph{validators} assert over the result. If the model errors, the command fails
to execute, or a validator rejects the outcome, the reflexer does not improvise:
it \emph{escalates} to the thinker with a tagged reason. Success is recorded; the
loop ends.

\begin{figure}[t]
  \centering
  \includegraphics[width=0.98\textwidth]{fig_reflexer.pdf}
  \caption{Inside the reflexer. The pattern's command shape and slot descriptions
  prompt the 2B model, which emits a single shell command in one forward pass; the
  command is executed and validated, and any failure escalates to the thinker
  rather than being retried by reasoning.}
  \label{fig:reflexer}
\end{figure}

A single forward pass is the entire reflexer runtime. There is no tool-use loop,
no planning step. When routing is correct and a pattern exists, this
constrained executor reaches \textbf{100\% binding accuracy} on the atom task set
(Chapter~\ref{ch:eval}). The point worth stating early is structural rather than
numerical: the reflexer succeeds not because the 2B model is strong, but because
the pattern has removed almost everything it could fail at.

\section{The thinker: a general agent for the tail}
\label{sec:thinker}

The thinker is a full agent loop in the ReAct style~\citep{yao2022react}: the
model is given the relevant skill documentation and a set of tools (a shell, a
file reader), and it reasons, acts, observes, and repeats until it judges the
task complete. It is invoked for everything the reflexer cannot serve:
multi-intent or compound requests, requests requiring clarification or planning,
and any intent for which no pattern has been distilled.

Concretely, the thinker is built on the \textbf{Pi coding-agent
framework}~\citep{pi-agent}, a minimal and self-contained agent loop. Pi was
chosen deliberately: among available agent frameworks it is the simplest and most
functional---a small, transparent loop with a clean tool interface and no heavy
orchestration layer---which keeps the thinker easy to reason about and to
instrument for the binding measurements of Chapter~\ref{ch:eval}. The same
property makes Pi the intended home for future work: the planned extensions of the
reflexer (Chapter~\ref{ch:conclusion}) are to be developed inside the Pi
environment, so that the executor and the agent loop share one runtime and tool
surface.

The thinker is where capability lives, but it is also where the architecture
pays its costs. An agent loop emits many commands across several turns, consumes
substantially more tokens and wall-clock time than a single reflexer call, and
at the parameter scale required for reliable reasoning it must run off-device.
A recurring theme of the evaluation is that raw thinker capability does not
translate cleanly into reliable \emph{binding}: a free-form model, even a large
one, tends to substitute its own pretraining priors (real Linux commands,
documented backend mechanisms) for the skill's defined vocabulary. The thinker
is therefore the right tool for novel problems and the wrong tool for the
repeatable head---precisely the asymmetry the two-tier split exploits.

\section{Patterns: atoms and composites}
\label{sec:patterns}

A \textbf{pattern} is the artifact that makes the reflexer possible. It is a
small, immutable JSON record that encodes one repeatable intent. Each pattern
carries: a trigger (how the router recognizes the intent), a set of typed
\emph{slots} (the variable parts to extract from the request), an
\texttt{args\_template} (the exact command shape, with slot placeholders),
optional validators (constraints on slot values), and natural-language example
queries used by the router's embedding index.

Patterns come in two kinds. An \textbf{atom} encodes a single action---set the
screen brightness, start a timer, toggle Wi-Fi---and binds to exactly one
command. A \textbf{composite} encodes a routine made of several atoms executed in
sequence, such as a ``good night'' action that dims the screen, enables Do Not
Disturb, and sets an alarm. Atoms are the unit of reliability: they are stable,
reusable, and in our measurements essentially perfectly reflexable. Composites
are more demanding, because some require values produced at runtime (a contact
name, a URL) rather than known in advance; these remain the harder case for the
reflexer and, where they need open-ended slot extraction, are better left to the
thinker.

Critically, the \texttt{args\_template} is a \emph{hard} constraint, not a
suggestion. Where the thinker reads skill prose and may reason its way to a
different command, the reflexer is handed the command shape directly and only
fills the slots. This distinction---soft prose for the thinker versus a hard
template for the reflexer---explains much of the reliability gap between the two
tiers and is revisited in Chapter~\ref{ch:eval}.

\section{The skill model: deep modules behind a CLI}
\label{sec:skills}

A \textbf{skill} packages a device capability behind a command-line interface
and a documentation file (\texttt{SKILL.md}). The CLI is the only surface either
tier touches: the reflexer fills a template that calls it, and the thinker reads
the \texttt{SKILL.md} and invokes it through the shell. Skills emit structured
results (a JSON object reporting success, data, and errors), so that outcomes can
be checked uniformly regardless of which tier produced the command.

The skill model is an application of Ousterhout's principle of \emph{deep
modules}: a simple interface that hides substantial implementation
complexity~\citep{ousterhout2018philosophy}. A well-designed skill exposes a
small, intent-shaped vocabulary (\texttt{clock timer --duration 10m}) while
absorbing the messy device-specific mechanism behind it. This depth is what lets
even a small model act competently when it is given the right tool: on a clean,
deep settings interface the 2B reflexer succeeds on multi-step tasks that it
fails when forced to manipulate the raw state directly (Chapter~\ref{ch:eval}).
The corollary, also measured, is that a \emph{leaky} skill---one whose
documentation exposes its backend mechanism---degrades the thinker, which
reasons toward the familiar mechanism instead of the intended alias. Interface
design is therefore not cosmetic; it is a determinant of binding accuracy.

\section{Distillation: from teacher to a symbolic pattern}
\label{sec:distillation}

Patterns are not written by hand at scale; they are \emph{distilled}. A capable
\textbf{teacher} model is used once, offline, to convert a skill's documentation
and device conventions into the pattern that the cheap on-device
\textbf{student} (the reflexer) then applies forever. This is the mechanism that
pays for the reflexer's narrowness: the one-time cost of minting a pattern is
amortized over every subsequent on-device invocation.

\begin{figure}[t]
  \centering
  \includegraphics[width=0.95\textwidth]{fig_distillation.pdf}
  \caption{Symbolic distillation. A strong teacher reads the skill documentation
  and schema and mints a frozen JSON pattern once, offline; the cheap 2B student
  applies that pattern on-device on every subsequent request. The teacher's cost
  is paid at distillation time, not at inference time.}
  \label{fig:pipeline}
\end{figure}

The relationship to classical knowledge distillation is deliberate and worth
making precise. In the original formulation, a large teacher network is
compressed into the \emph{weights} of a smaller student
network~\citep{hinton2015distilling}. A.L.F.R.E.D.\ distills along the same
teacher$\rightarrow$student axis, but into a different target: a teacher model
compresses its knowledge of a device's conventions into an \emph{inspectable
symbolic artifact}---the pattern---rather than into student parameters. The trade
is explicit. Weight distillation yields a general student but requires
retraining to extend and is opaque to inspection. Symbolic distillation yields a
student that is plug-and-play (adding a capability is adding a JSON file, with no
retraining), auditable (the encoded convention can be read and corrected), and
composable, at the cost of covering only the intents that have been distilled.
For an on-device system that must be extended by skill authors and trusted by
users, that trade is the right one.

Two regimes of distillation appear in this work. In \textbf{cold-start}
distillation the teacher sees only the skill documentation and the API schema,
never any task or its expected outcome; it must encode each convention from
documentation and prior knowledge alone. In \textbf{session-based} (grounded)
distillation the teacher additionally observes the real device---reading back
actual setting codes before committing an encoding. Cold-start is cheaper and
captures well-documented conventions, but it cannot recover genuinely
device-specific values that no amount of reasoning can derive; those require
grounding. The boundary between what cold-start captures and what only grounding
can fix is one of the empirical results of Chapter~\ref{ch:eval}. A strict
integrity rule governs both regimes: the teacher must never see the evaluation
tasks or their gold outcomes, and a distilled pattern, once minted, is frozen
and never edited after observing failures---editing it would leak the benchmark
into the artifact under test.

\section{Design principles}
\label{sec:principles}

Three principles run through the design and are stated here because they recur as
justifications later.

\paragraph{Contract-first.} Each benchmark is specified before it is built. The
invariants (state isolation, a single success oracle, the metrics) and the
oracle itself are written down as a contract, and only then is the mechanism
implemented against that contract. This ordering prevents the common failure of
letting an evaluation drift to match whatever the system happens to produce.

\paragraph{Success is an end state, not a command.} A.L.F.R.E.D.\ is judged on
whether the device reaches the requested state without collateral changes, not on
whether it emitted a particular command string. This catches two failures that
command matching misses: a syntactically valid command that produces the wrong
outcome, and an over-eager agent that achieves the goal but also changes things
it was not asked to. The end-state oracle that enforces this is detailed in
Chapter~\ref{ch:method}.

\paragraph{Distillation is the learning phase; the reflex is frozen.} All
adaptation, iteration, and feedback live in the offline teaching step. The
runtime reflex is a single pass over a fixed pattern and never changes during
execution. This separation keeps the on-device path simple and predictable
however much effort the teacher spends producing the pattern---the analogue of a
practiced human reflex, refined slowly through feedback yet executed instantly
and without deliberation.

\bigskip
\noindent
Together these pieces define the system under test: a router that splits traffic,
a reflexer that executes distilled patterns in one shot, a thinker that handles
the tail, skills that present deep interfaces, and a distillation step that mints
the patterns offline. The chapters that follow put numbers to every claim made
here---how much model binding really needs, when the abstraction helps, where
scale stops helping, and how much the choice of teacher matters.
